\documentclass[11pt]{article}
\usepackage{fontspec}
\usepackage[french]{babel}
\usepackage{amsmath}
\usepackage{amssymb}
\usepackage{amsthm}
\usepackage{mathrsfs}
\usepackage{tikz-cd}
\usepackage{enumitem}
\usetikzlibrary{babel}

\newcommand{\printedpage}[1]{\marginpar{\footnotesize p.~#1}}
\newcommand{\msp}[1]{\marginpar{\footnotesize [#1]}}
\newcommand{\Cathot}{\mathrm{Cat\,hot}}
\newcommand{\uHom}{\underline{\mathrm{Hom}}}
\newcommand{\uuK}{\underline{\underline{K}}}
\newcommand{\uCh}{\underline{\mathrm{Ch}}}
\newcommand{\Hom}{\mathrm{Hom}}
\newcommand{\Hot}{\mathrm{Hot}}
\newcommand{\Ens}{\mathrm{Ens}}
\newcommand{\Ch}{\mathrm{Ch}}
\newcommand{\defeq}{\overset{\text{déf}}{=}}
\newcommand{\Chh}{\underline{\underline{\mathrm{Ch}}}}

\begin{document}

\printedpage{29}
% head: A. Grothendieck : Dérivateurs, Ch. V, Édition des Universités de Montpellier II et Paris VII
% foot: Version 9 septembre 2001
$K_\infty$, dont la réunion est $K_\infty$. Mais d'autre part je ne sais donner aucune interprétation convaincante de $K_\infty$, comme de celle du ind-objet de Cat\,hot defini par $\uuK = (K_n)_{n\in\mathbf{N}}$, qui ind-représente le foncteur remarquable % [sic] defini -> défini, missing accent
\[
H = H_{\mathcal{M};x,y} : \Cathot^{o} \longrightarrow \Ens
\]
factorisant le foncteur
\msp{page 45}
\[
X \longmapsto H_{\mathcal{M};x,y}(X) \defeq \Hom_{\mathcal{M}(X)\Sigma_X^{-1}}(x_X, y_X),
\]
par la formule donc
\[
H(X) = \varinjlim_{n} \underbrace{\pi_0\uHom(X, K_n)}_{\Hom_{\Cathot}(X,K_n)} .
\]
Je ne sais donner d'interprétation de ce type ni de la catégorie $K_\infty$ elle-même, ni de son image dans Cat\,hot, ni de son image dans Hot. Il est vrai que si je restreins le foncteur $H$ à la sous-catégorie pleine de (Cat) formée des catégories \textit{finies} (Ob\,$X$ et Fl\,$X$ finis), alors on aura
\[
\uHom^{!}(X, K_\infty) = \uHom(X, K_\infty),
\]
et l'isomorphisme précédant s'interprète comme % [sic] précédant -> précédent, adjective intended
\[
\Hom_{\mathcal{M}(X)\Sigma_X^{-1}}(a_X, b_Y) \defeq H(X) \simeq \pi_0\uHom(X, K_\infty) \defeq \Hom_{\Cathot}(X, K_\infty).
\]
\texttt{[plutôt} $(a_X, b_X)$\texttt{]} Cela prouve du moins que la connaissance de $K_\infty$ est utile. Mais cette formule, pour $X$ restreint aux catégories \textit{finies,} ne définira pas $K_\infty$ comme objet de Cat\,hot, à isomorphisme unique près dans cette catégorie - \textit{sauf} dans le cas où $K_\infty$ est elle-même homotope à une catégorie
\msp{page 46}
finie. (Je ne me rends pas bien compte si ça arrive très souvent \ldots\ Mais on ne peut fonder une théorie générale des dérivateurs et des catégories de modèles `dérivables' sur une hypothèse de ce genre.)

Mais supposons que le foncteur $H$ envisagé se factorise même par Hot, i.e. qu'il transforme quasi-isomorphismes en bijections. On a vu que cela équivaut à la condition qu'il puisse se représenter par un système inductif sur $\mathbf{N}$ formé de catégories de Thomason-Kan. Il doit être facile de plus, si on y tient, de s'arranger pour que les foncteurs de transition soient des monomorphismes et pleinement fidèles, pour avoir une $\varinjlim$ sympathique. Ce sera bien sûr une catégorie de Thomason-Kan. Cette catégorie est-elle déterminée par le foncteur $H : \Cathot^{o} \longrightarrow \Ens$ à isomorphisme unique près dans Cat\,hot? Bien sûr, il n'est pas question qu'elle représente le foncteur $H$, sauf \texttt{[si]} celui-ci est bel et bien représentable, i.e. si le ind-objet $\uuK$ de Cat\,hot est essentiellement constant (ce qui ne doit pratiquement jamais être le cas). Notons déjà
\printedpage{30}
% head: A. Grothendieck : Dérivateurs, Ch. V, Édition des Universités de Montpellier II et Paris VII
% foot: Version 9 septembre 2001
\msp{page 47}
que si $u$ est un homomorphimse de ind-objets \textit{dans} Cat, ou, ce qui revient maintenant au même, un homomorphisme dans Cat % [sic] homomorphimse -> homomorphisme, transposed letters
\[
K_\infty \xrightarrow{u_\infty} K'_\infty
\]
tel que $\forall n$, $u_\infty(K_n) \subseteq K'_{n'}$ pour $n' = \tau(n)$ convenable, et si on suppose que l'homomorphisme de ind-objets induit un isomorphisme pour les foncteurs ind-représentables \textit{sur} Cat\,hot (par exemple, si $\uuK$ et $\uuK'$ ind-représentent le même foncteur $H$, et l'homomorphisme de ind-objets est celui qui correspond à $\mathrm{id}_H$), alors $u_\infty$ est forcément un isomorphisme \textit{dans} Hot, donc aussi (sauf erreur) dans Cat\,hot puisque $K_\infty$, $K'_\infty$ sont des catégories de Thomason-Kan. Mais ce qui n'est pas évident, c'est l'existence d'un homomorphisme $u_\infty$ comme \texttt{[ci-]}dessus, et s'il existe, qu'il soit unique (ou du moins canonique?) comme flèche de Cat\,hot.

Pour l'existence d'un $u_\infty$, on est ramené à prouver que si on a un diagramme à flèches pleines
\[
\begin{tikzcd}
K_n \arrow[r, "i_n", "\text{cofib.}"'] \arrow[d, "u_n"] & K_{n+1} \arrow[dr, dashed] & \\
K'_{\tau(n)} \arrow[rr, "\text{transition}"'] & & K'_N ,
\end{tikzcd}
\]
et si une flèche $u_{n+1}$ est donnée `en pointillés' \textit{dans} Cat\,hot, de façon à rendre le carré commutatif,
\msp{page 48}
alors on peut la réaliser par une flèche \textit{dans} Cat qui rende le diagramme commutatif pour de bon (dans Cat). On est ramené au même problème dans une situation
\[
\begin{tikzcd}
X \arrow[r, hook, shift right, "i"'] & Y , \arrow[l, dashed, shift right, "p"']
\end{tikzcd}
\]
où $X$ est une catégorie de Thomason, $i$ une cofibration de Thomason, $p$ une rétraction de $Y$ dans $X$ \textit{dans} Cat\,hot, on veut la réaliser par une vraie rétraction (\footnote[7]{\texttt{[un ?\ \ dans la marge]}}).

En admettant que ça soit le cas, donc que tout homomorphisme des ind-objets $\uuK$, $\uuK'$ \textit{dans} Cat\,hot se représente par un homomorphisme entre les ind-objets \textit{dans} Cat, donc aussi par un foncteur $K_\infty \longrightarrow K'_\infty$, reste la question si l'homomorphisme correspondant dans Cat\,hot est unique. Donc si on a deux foncteurs $f$, $g$ qui induisent le même homomorphisme sur les foncteurs ind-représentés sur Cat\,hot, si $f$ et $g$ sont homotopes. Donc

\textbf{Question}. Deux foncteurs $f, g : K_\infty \longrightarrow K'_\infty$ tels que pour tout $n \in \mathbf{N}$, $f_n$, $g_n$ qui s'en déduisent par restriction à $K_n$ sont homotopes (et même homotopes en tant que foncteurs de $K_n$ dans un $K'_N$ convenable), sont-ils homotopes? (\footnote[8]{$K$, $K'$ catégories de Thomason, les $K_n$, $K'_n$ aussi.})

\printedpage{31}
% head: A. Grothendieck : Dérivateurs, Ch. V, Édition des Universités de Montpellier II et Paris VII
% foot: Version 9 septembre 2001
\msp{page 49}
Je présume que non - qu'il peut arriver que la longueur minimum requise $l(n)$ pour un chemin dans $\uHom(K_n, K'_{\tau(n)})$ reliant $f_n$ et $g_n$, tende vers l'infini si $n \longrightarrow \infty$. L'hypothèse que les $K'_n$ sont des catégories de Thomason y change-t-il quelque chose (deux objets homotopes dans $\uHom(X, K')$, si $K'$ est de Thomason, sont-ils reliés par une flèche?)

Bien sûr, on peut prétendre que les catégories $\uCh_{\Sigma,n}(\mathcal{M}; a, b)$ ont été définies `canoniquement', et leur limite inductive $\uCh_{\Sigma}(\mathcal{M}; a, b)$ aussi. Mais la définition était assez arbitraire et abracadabrante - je suis convaincu que dix mathématiciens, se proposant de développer cette même idée des catégories de chemins, dans le même but que moi, arriveraient à dix définitions différentes! Pourtant, il n'est pas exclu que l'on finisse par trouver un système transitif d'isomorphismes dans Cat\,hot, entre les catégories de chemins `infinis' (i.e. de longueur non précisée) $K_\infty = K_\infty(\mathcal{M}; a, b)$ que nous aurons construites, dans le but
\msp{page 50}
principal (disons) de représenter le foncteur $H$ \textit{sur les catégories finies.} Et que les isomorphismes obtenus seront, bien entendu, des homomorphismes dans Ind(Cat\,hot). J'ai l'impression de ne pas avoir encore axiomatisé à fond la notion même de chemins, qui après tout a un sens géométrique assez évident et important par lui-même. Donc par ce qu'il faut entendre par `\textit{une} théorie des chemins dans une catégorie'. Une telle théorie $\mathcal{T}$ associerait donc à $\mathcal{M}$, $\Sigma$, $a$, $b$ un système inductif de catégories $\Ch^{\mathcal{T}}_{n,\Sigma}(\mathcal{M}; a, b)$. Et les contraintes géométriques naturelles à imposer à de telles théories, en vue aussi de leur utilisation pour expliciter
\[
\Hom_{\mathcal{M}(X)\Sigma_X^{-1}}(a_X, b_X) \xrightarrow{\ \sim\ } \pi_0\uHom(X, \Chh^{\mathcal{T}}(\mathcal{M}; a, b)) % reading resolved: two underlines, the ind-path object.
\]
donnerait un système transitif d'isomorphismes non seulement entre les ind-objets de Cat\,hot, mais peut-être aussi entre les $\Ch^{\mathcal{T}}_{\infty,\Sigma}(\mathcal{M}; a, b)$, pour des théories differentes. Si on dispose d'un tel énoncé, la catégorie $\Ch_{\infty,\Sigma}(\mathcal{M}; a, b)$, ou plus exactement ses images dans Cat\,hot et Hot, % [sic] differentes -> différentes, missing accent
\msp{page 51}
prendraient un sens convaincant.

Il y a d'autre part un autre point de vue, qui peut-être pourrait `canonifier' les catégories $\uCh_\infty$ en tant qu'objet de Cat\,hot et surtout de Hot (on passerait à Cat\,hot en appliquant le foncteur $\psi$, enveloppe de Thomason-Kan \ldots). En effet, on est intéressé à construire, en termes de $\mathcal{M}$, $\Sigma$, un `dérivateur' (ou un prédérivateur) $\mathbf{D}_{\mathcal{M},\Sigma}$ surtout sur le `domaine' des catégories finies. L'étendre à un domaine plus large est peut-être important dans un `deuxième souffle', mais peut être considéré comme du luxe dans le premier mouvement de fondation d'une théorie des dérivateurs. Ceci dit, on peut se poser des questions du type suivant.

Soit
\[
\Hot_0 \subset \Hot
\]
\printedpage{32}
% head: A. Grothendieck : Dérivateurs, Ch. V, Édition des Universités de Montpellier II et Paris VII
% foot: Version 9 septembre 2001
la sous-catégorie pleine de Hot, définie par les catégories finies (\footnote[9]{\textit{NB} Hot$_0$ est équivalente à une petite catégorie.}) - appelons les types d'homotopie `de type fini' (ou de \textit{présentation} finie?). On a un foncteur évident
\msp{page 52}
\[
\begin{aligned}
\Hot \quad &\longrightarrow \quad \Hot_0^{\wedge} = \uHom(\Hot_0^{o}, (\Ens)) \, ,\\
X \quad &\longmapsto \quad \tilde{X} = (\xi \longmapsto \Hom_{\Hot}(\xi, X)).
\end{aligned}
\]
Ce foncteur est \textit{conservatif} \texttt{[cela semble faux]}, puisque un homomorphisme $X \longrightarrow X'$ dans Hot est un isomorphisme s'il induit des isomorphismes pour les foncteurs $\tilde{X} \longrightarrow \tilde{X}'$, et même seulement pour leurs valeurs pour les sphères $S^n$ ($n \geq 0$). (Sauf erreur. Y a-t-il des canulars, parce que le critère d'équivalence faible par les $\pi_i$ ne s'applique pas directement, en l'absence de points bases. Ça ne doit pas être très sérieux \ldots\ \texttt{[cela semble en fait sérieux]}). Mais est-il peut-être même \textit{pleinement fidèle} (\footnote[10]{\texttt{[un ?\ \ dans la marge]}})? De sorte que Hot se réaliserait comme sous-catégorie pleine de $\Hot_0^{\wedge}$, contenant les foncteurs représentables?

\end{document}
